\item Un sifón se utiliza para drenar agua de un tanque, como se ilustra en la figura P14.31. Suponga flujo estable sin fricción. (a) Si $h = 1.00 \text{ m}$, encuentre la rapidez del flujo de salida en el extremo del sifón. (b) ¿Qué pasaría si? ¿Cuál es la limitación en la altura de la parte superior del sifón respecto al extremo de éste? Nota: Para que el flujo del líquido sea continuo, su presión no debe caer debajo de su presión de vapor. Suponga que el agua está a $20.0^\circ\text{C}$, con una presión de vapor de $2.3 \text{ kPa}$.
